\subsection{Eventual publication and the source-selection problem}
\label{sec:source-eventual-publication}

A finite deadline can be eliminated from the declared native record model.
This gives compatible temporal laws and, on paths with separated source
coordinates, an unconditional rank invariant for every finite native word.
The result does not identify the public meaning, population or radius with
those selected by the complete A1--A3 source architecture.

\paragraph{Constructed continuation potential.}
Let $T_a s$ advance a response checkpoint by a permitted native move and
let $p(s,a)>0$, $\sum_a p(s,a)=1$, be its declared finite-alphabet reference.
Write $f(s)$ for the indicator that retained observations determine the
required public meaning. Retention makes publication persistent. The
exhaustive finite continuation masses satisfy
\[
 H_0=f,\qquad H_{n+1}(s)=\sum_a p(s,a)H_n(T_a s),\qquad
 h(s)=\sup_n H_n(s).
\]
They increase from $f$ to a limit in $[0,1]$. Finite summation commutes
with the limit, so $h(s)=\sum_a p(s,a)h(T_a s)$. Induction proves that
$h$ is the least harmonic majorant of $f$. Its positivity is equivalent
to a finite publishing continuation; the connected native completion
theorem identifies this with preservation of the required meaning by
the joint retained-record/current-state map. The initially sampled
receiver hypothesis is essential.

For $h(s)>0$, the cylinders and transition probabilities
\[
 Q_s[w]=\frac{R_s[w]h(T_w s)}{h(s)},\qquad
 K(s,a)=\frac{p(s,a)h(T_a s)}{h(s)}
\]
are normalized and obey $\sum_a Q_s[wa]=Q_s[w]$. This is the standard
Doob transform, with its potential constructed from native publication.\footnote{
For the standard transform see D.~A.~Levin and Y.~Peres, with E.~L.~Wilmer,
\emph{Markov Chains and Mixing Times}, second edition, section~17.6,
\url{https://pages.uoregon.edu/dlevin/MARKOV/markovmixing.pdf}.}
The actual length-$n$ prefix marginal of the finite deadline-$n+k$
information projection is
\[
 \frac{R_s[w]H_k(T_w s)}{H_{n+k}(s)}\longrightarrow Q_s[w].
\]
An unpublished but recoverable prefix has zero mass at its own deadline
when that deadline is feasible, and positive limiting mass. An infeasible
deadline has no conditioned probability law. This limiting construction resolves the
finite-deadline incompatibility without identifying time restriction with
regulator refinement.

Analytically, the finite-alphabet reference defines a probability on
infinite words. Publication at some finite time is a countable union of
cylinder events, of probability $h(s)$. Its conditional law has exactly
the displayed cylinders and publishes almost surely. For any competing
path law $V$ supported on this event and absolutely continuous with respect
to the reference,
\[
 D(V\Vert R)=D(V\Vert Q)-\log h(s).
\]
Singular laws have infinite cost. Hence $Q$ is the unique information
projection on this declared infinite history algebra. This is an analytic
extension of the finite schedule problem, not a replacement of canonical
A3's required source grammar by a chosen history algebra.

\paragraph{Quantitative completion.}
Suppose there are $N\geq2$ ports and all reference move probabilities are
at least $\eta>0$ at every reachable checkpoint. Connected completion
supplies a word of at most $B=N(N-1)/2$ means from every viable checkpoint.
Padding preserves publication, so its reference probability is at least
$\delta=\eta^B$. Published, dead and viable unpublished states give
\[
 h-H_B\leq(1-\delta)(h-f),\qquad
 0\leq h-H_{kB}\leq(1-\delta)^k(h-f).
\]
The second inequality follows by applying the positive transition operator
and iterating. Consequently the selected first-publication time satisfies
$Q(T>kB)\leq(1-\delta)^k$ and $\mathbb E_QT\leq B/\delta$.
For the effective result, assume rational preparation and meaning matrices
on a finite-dimensional real parameter space, an exact rational evaluator
for the reference probabilities and a supplied positive rational uniform
lower bound $\eta$. Finite enumeration and kernel elimination then compute
$H_n$ exactly, and
\[
 H_{kB}\leq h\leq
 \min\{1,H_{kB}/[1-(1-\delta)^k]\},\qquad k\geq1,
\]
give convergent rational enclosures. Refining these enclosures and random
bits implements exact transition sampling almost surely. No efficient
planning bound, fixed physical precision or physical random source follows.
Positivity alone does not supply effectiveness: the four-port reference
$((1-r)/2,r/2,1/2)$, for noncomputable $1/3<r<2/3$, has lower bound $1/6$
but publication potential $h=r$. The real-valued existence and tail results
do not imply its computability.
For $N=1$, viability and the initial receiver sample imply publication at
time zero.

\paragraph{Permanent native rank on a path.}
On the path $0,\ldots,d$, $d\geq2$, prepare independent unknown deviations
$x,y$ at $0,1$ and calibrated zeros elsewhere; retain samples at receiver
$d$. For the two response columns define the ordered minors
$\Delta_{ij}=x_i y_j-x_j y_i$, $i<j$. All are initially nonnegative.
An adjacent pair mean sends a minor containing both averaged rows to zero,
leaves one containing neither unchanged, and otherwise averages two ordered
minors. Nonnegativity is preserved. A fixed nonadjacent minor is therefore
at least half its previous value after every move.

If edge $0$ acts before edge $1$, it erases the independent difference
before any receiver sample. If edge $1$ acts first, then
$\Delta_{02}=1/2$ and after every $t$ subsequent native means
$\Delta_{02}\geq2^{-(t+1)}>0$. Both coordinates remain encoded in the
evolving scalar state forever. Connected completion and a uniform positive
reference bound give eventual publication with probability one on this
safe class. For constant probabilities $p_i>0$, the initial potential and
selected law are therefore
\[
 h=\frac{p_1}{p_0+p_1},\qquad K_0=0,\quad K_1=p_0+p_1,
 \quad K_i=p_i\ (i\geq2).
\]
After entry into the safe class the selected law is the reference. For
$d=1$ the receiver initially records $y$ and $h=1$ instead. Local deletion
and renormalization of the destructive move gives a different law: on
the uniform four-port path it is $(0,1/2,1/2)$, whereas the derived law is
$(0,2/3,1/3)$.

There is an analytic extension to $k$ independent preparation coordinates
at nonadjacent path vertices $I$. The initial $N$-by-$k$ response has
nonnegative maximal ordered minors and determinant one on $I$. Adjacent
pair-mean matrices are totally nonnegative. Cauchy--Binet preserves
nonnegativity; in the update of the $I$ minor, its own coefficient is one
or one half, since no adjacent pair lies in $I$. All other contributions
are nonnegative. Thus this minor is at least $2^{-t}$ after every finite
word of length $t$. The state retains full preparation rank, and every
prefix admits connected completion. Under uniformly positive references,
eventual publication of the full preparation has probability one:
$h=1$ and the publication constraint changes no path probabilities.
This removes the preservation guard and successful-word restriction in
this exact prepared path model. Extra non-path means and fixed physical
noise tolerance are outside the theorem; the determinant bound allows
exponential attenuation.

\paragraph{Selecting the demand on one common algebra.}
Requiring only $x+y$ on the same path gives $h=1$: native means conserve
that total and it remains viable from every prefix. Include both this
aggregate demand and the two-record demand as labels $m$ in a single
declared joint history algebra, with faithful prior $\pi_m$. Constrain
publication of the labelled meaning. The unique joint information
projection has marginal
\[
 Q(m)=\frac{\pi_m h_m}{\sum_j\pi_j h_j}.
\]
Every recoverable demand retains positive mass. For equal priors in the
uniform four-port control the two-record and aggregate probabilities are
$1/3$ and $2/3$. Strict convexity on this common algebra therefore does
not force the stronger demand. The complete A1-generated grammar must
identify the required meaning or exclude the competitors; the declared
label algebra is not asserted to be that grammar.
For separated source coordinates the conclusion is stronger: every demand
that is a function of the full preparation has $h_m=1$, hence $Q(m)=\pi_m$.
In particular, competing metric-neighbourhood read subsets keep exactly
their prior weights at finite population. Publication supplies no preference
among them in this native model.

\paragraph{Population and radius do not follow from the limit tests.}
For the golden sites, every fixed $0<\alpha<1$ gives radii
$a_q=Lq^{-\alpha}$ and durations $\Delta_q=a_q/c$ with
\[
 a_q\longrightarrow0,\qquad
 \frac{2\sqrt3L/q}{a_q}=2\sqrt3q^{\alpha-1}\longrightarrow0.
\]
The existing covering-cone theorem and continuity-set counting theorem
therefore apply throughout this family, under their stated boundary
hypotheses. The grid $Lb/q$, $0\leq b_i<q$, has equal-volume cells and
assignment error at most $\sqrt3L/q$, giving the same limiting tests.
The finite read menus differ. At $q=21$, for exponents $1/4,1/2,3/4$,
the golden ordered read counts including self are respectively
$20376393$, $2911845$, $327151$; the grid counts are
$20141913$, $3118389$, $277255$.

\paragraph{The square root as an error balance.}
The proved generated-volume estimate gives a concrete reason to use a
square-root scale. In units with $c=1$, on a window $0\leq T\leq1$, let
both covering and cell-assignment error be bounded by $0<h\leq1$ and
let $2h<a\leq1$. Under the existing partition and buffered-domain
hypotheses, the actual interval-volume error obeys
\begin{align*}
 E&\leq4\pi(T+a)(T/2+h)^2(h+Th/a)+\pi T^3a/2\\
  &\leq18\pi h+\pi(a/2+18h/a).
\end{align*}
The time quadrature contributes the $a$ term, while accumulated spatial
error contributes $h/a$. The radius-dependent bound has its unique
minimum at $a=6\sqrt h$, because
\[
 a/2+18h/a-6\sqrt h=(a-6\sqrt h)^2/(2a)\geq0.
\]
For $h<1/36$ the minimizer is admissible. With $h=C/q$ its exponent is
one half, but its coefficient is $6\sqrt C$, not the prescribed unit
coefficient. This optimizes a certified upper bound, not necessarily the
actual error or canonical A3's entropy. Restoring a reference length gives
the geometric-mean form $\sqrt{Lh}$. The exact radius is therefore
dispensable for the stated continuum target and has a defensible role as
a regulator; it has not been derived as a physical source length.

\paragraph{Removing locality fails; a fixed stencil also need not suffice.}
If every event reads every site in the preceding layer, the existing
generated read relation satisfies
\[
 e\preceq f\quad\Longleftrightarrow\quad
 \operatorname{layer}(e)<\operatorname{layer}(f)\ \text{or}\ e=f.
\]
A fixed positive spatial separation is crossed in one vanishing layer
duration, violating every fixed finite speed. An interval with two
singleton tips and $K$ full interior layers of $N$ sites has ordering
fraction
\[
 F=1-\frac{KN(N-1)}{(KN+2)(KN+1)},\qquad 1-1/K\leq F\leq1.
\]
It tends to one uniformly in $N\geq1$, not the four-dimensional
Minkowski interval value $1/10$. The count includes both tips and compares
actual intervals, not rectangular windows.

Nor does an arbitrary microscopic local stencil force a round cone.
On the unit cube grid with spacing $1/q$, radius $a=1/q$ and duration
$a$, only axial neighbours and self are read. Reaching $(1/2,1/2,0)$
from zero takes time one, although it is strictly Euclidean-timelike
at time $3/4$. This persists for every $q$ divisible by four: the cone
uses the $\ell^1$ norm. At radius $1/q^2$ only self reads remain.
These controls explain the separated scales of the successful family;
they do not rule out every radius-free architecture. A source-derived
read relation could replace the ball prescription if its actual
propagation and counting limits were proved. The exact finite M1 menu
and that continuum target are distinct obligations. None of these
geometric controls is a full A1--A3 countermodel.

\paragraph{The full finite A3 state problem.}
Let $\mathcal K$ be A3's nonempty convex family of compatible local states,
with finite observer cover $\mathcal G$, weights $w_P>0$, faithful local
references $\tau_P$, and attained objective
\[
 F(\rho)=\sum_{P\in\mathcal G}w_P D(\rho_P\Vert\tau_P).
\]
The algebras may be noncommutative, and the local family need not extend
to a global state. Every minimizer $\rho_*$ satisfies
\begin{equation}
 \operatorname{supp}\sigma_P\subseteq\operatorname{supp}(\rho_*)_P
 \quad(\sigma\in\mathcal K,\ P\in\mathcal G).
 \label{eq:source-state-support}
\end{equation}
Thus for every positive effect $E$ in a scored algebra,
\begin{equation}
 (\rho_*)_P(E)=0
 \quad\Longleftrightarrow\quad
 \sigma_P(E)=0\ \hbox{for every }\sigma\in\mathcal K.
 \label{eq:source-state-zero}
\end{equation}
An actual positive restriction pullback transfers this result to a public
record; injectivity of a cover alone does not provide that pullback.

For $\rho_t=(1-t)\rho+t\sigma$, $0<t<1$, trace expansion gives
\begin{align*}
 &(1-t)D(\rho_P\Vert\tau_P)+tD(\sigma_P\Vert\tau_P)
          -D((\rho_t)_P\Vert\tau_P)\\
 &\qquad=(1-t)D(\rho_P\Vert(\rho_t)_P)
              +tD(\sigma_P\Vert(\rho_t)_P).
\end{align*}
The logarithms on the right are restricted to the mixture support,
which contains both endpoint supports. If $Q$ projects onto the kernel of
$\rho_P$ and $b=\operatorname{Tr}(Q\sigma_P)>0$, binary measurement data
processing bounds this gap below by $-tb\log t-tH(b)$, with
$H(b)=-b\log b-(1-b)\log(1-b)$. All other local gaps are nonnegative.
Consequently
\[
 F(\rho_t)-F(\rho)
 \leq t\bigl[F(\sigma)-F(\rho)+w_PH(b)\bigr]+tw_Pb\log t.
\]
At a minimum the left side is nonnegative, whereas the right side is
negative for small enough $t$. This proves \eqref{eq:source-state-support}
without simultaneous diagonalization. Zero expectation of a positive
effect means it annihilates the state support, proving
\eqref{eq:source-state-zero}. The imported data-processing theorem is
Theorem~5.35 of Watrous, \emph{The Theory of Quantum Information},
\url{https://cs.uwaterloo.ca/~watrous/TQI/TQI.pdf}.

On the established common supports, the derivative at $t=0$ gives
\[
 \sum_P w_P D(\sigma_P\Vert(\rho_*)_P)
       \leq F(\sigma)-F(\rho_*).
\]
For $b=\sigma_P(E)>0$, binary data processing and
$d(b\Vert p)\geq-H(b)-b\log p$ therefore imply
\begin{equation}
 (\rho_*)_P(E)\geq
 \exp\!\left[-\frac{F(\sigma)/w_P+H(b)}{b}\right]>0.
 \label{eq:source-state-failure-floor}
\end{equation}
This uses a feasible witness, without solving the optimizer. It is not
uniform across refinements without uniform witness-cost and weight bounds,
so it does not exclude an asymptotically vanishing error. For a candidate
with zero effect probability, set $a=w_Pb$, $C=F(\sigma)+w_PH(b)$.
The mixture $t=\exp(-1-C/a)$ improves its objective by at least
$a\exp(-1-C/a)$; an approximation tolerance below this value cannot
justify the candidate. For singular references, compress to their
supports and apply the result to the nonempty finite-objective face of
$\mathcal K$. A zero in a chosen reference can exclude an otherwise
feasible failure, but its source derivation then requires proof.

\paragraph{A selected state does not determine its record process.}
Exact noncommuting controls use $R=|0\rangle\langle0|$,
$S=|v\rangle\langle v|$, rational unit vectors
$v=(3/5,4/5),(5/13,12/13),(8/17,15/17)$, and reference $I/2$.
The unitary $\left(\begin{smallmatrix}a&b\\b&-a\end{smallmatrix}\right)$
exchanges the endpoints, so strict convexity gives the unique midpoint
optimizer, with determinant $b^2/4>0$ and last-coordinate probability
$b^2/2$, although that probability is zero at $R$.

Identity and complete depolarization both preserve $I/2$ and are
covariant under every unitary, but the mismatch probabilities of their
uniform classical two-time records are zero and $1/2$. On the declared
common set of normalized channel Choi states with input marginal $I/2$,
reference $I/4$ selects depolarization. The bit-flip channel is a
deterministic-failure witness of cost $\log4$, giving the lower bound
$1/4$ in \eqref{eq:source-state-failure-floor}. Adding the zero-mismatch
face instead selects $\operatorname{diag}(1/2,0,0,1/2)$, the dephasing
channel. This preserves classical records while discarding phase; it is
not the identity quantum channel. This process state space and constraint
are inputs, not a reconstruction of A1's instruments or a full-axiom
countermodel. The complete proofs are in
\path{code/source_publication_selection/STATE_SELECTION.md}.

For an identified positive failure effect, \eqref{eq:source-state-zero}
settles the finite A3 selection step: it enforces zero failure precisely
when the source feasible family excludes failure. Establishing
that family and the meaning of its record effects remains necessary.
Conditioning on desired success cannot derive those inputs.

\paragraph{Proof and source boundaries.}
The modules under \path{Lean/Geometry/SourcePublicationAxiomAudit.lean}
check the continuation potential, compatible cylinders, finite-marginal
limit, native viability, block contraction, all-word two-column invariant,
joint-demand marginal, radius/count formulas, complete-read order,
interval-fraction bounds and the volume-error balance. Infinite-path measure and
entropy arguments, the uniform native block instantiation and general
$k$-column Cauchy--Binet argument are analytic. The independent controls
in \path{code/source_publication_selection/} replay complete finite word
trees and all maximal minors, and certify golden/grid metric counts by
different exact arithmetic methods. Further controls independently
reconstruct complete-read interval closure and exact grid-stencil paths.
The noncommutative support, derivative and quantitative entropy arguments
are analytic; \path{SourceStateSelection.lean} kernel-checks their real
inequality reductions, with the matrix gap estimate as an explicit
hypothesis. It does not formalize Umegaki data processing. Exact matrix
and channel replay supplies independent regression controls.
A full A1--A3 extraction of the golden
population, required immutable-version meaning and particular radius,
compatible with actual source response, transport and regulator refinement,
remains necessary for literal finite M1. The continuum causal target can
dispense with that exact radius, but requires a source-derived local
propagation law. The results here do not discharge that obligation.
